\origpage{376}{N. Robertson, P.D. Seymour / Journal of Combinatorial Theory, Series B 90 (2004) 325--385}
Since each $(\Gamma, \phi, \lambda, \mathscr{T}) \in \mathscr{D}'$ has the same $\omega_0$-flaw, it follows that there is a $\mathscr{W}$ such that $\mathscr{W}$ is a cover for each $(\Gamma, \phi, \lambda, \mathscr{T}) \in \mathscr{D}'$; and for each $W \in \mathscr{W}$ there exists $H(W)$ and $\rho(W)$ such that $H(W)$ is a $W$-hub of radius $\rho(W)$ for each $(\Gamma, \phi, \lambda, \mathscr{T}) \in \mathscr{D}'$ (they have the same radius in each $\Gamma$ because the common $\omega_0$-flaw includes a spoke). Moreover, for each $W \in \mathscr{W}$ the sets $S(W), N(W)$ are the same for each $(\Gamma, \phi, \lambda, \mathscr{T}) \in \mathscr{D}'$, and the spokes $S_v$ $(v \in N(W))$ are also the same for each $(\Gamma, \phi, \lambda, \mathscr{T})$.

Thus we may speak of $S(W), N(W), S_v$ $(v \in N(W))$ and $\Delta_v$ $(v \in S(W))$ without specifying a particular member of $\mathscr{D}'$. (This does not hold for general spokes, of course; they will differ for different members of $\mathscr{D}'$, as indeed will the sets $V(\Gamma) \cap bd(W)$.) For each $W \in \mathscr{W}$, fix a march $\pi(H(W))$ with $\bar{\pi}(H(W)) = V(\Gamma) \cap V(H(W))$ for each $(\Gamma, \phi, \lambda, \mathscr{T}) \in \mathscr{D}'$, and for each $r^* \in S(W)$ fix a march $\pi(r^*)$ with $\bar{\pi}(r^*) = V(\Gamma) \cap V(S_{v_1} \cup S_{v_2})$ for each $(\Gamma, \phi, \lambda, \mathscr{T}) \in \mathscr{D}'$, where $v_1, v_2$ are the neighbours of $r^*$ in $\Gamma^* \cap bd(W)$ (and hence $v_1, v_2 \in N(W)$).

Let $\Sigma'$ be obtained from $\Sigma_\chi$ by deleting $\Sigma_\chi \cap (W \setminus bd(W))$ for each $W \in \mathscr{W}$. Then $\Sigma'$ is a surface since $bd(W) \cap bd(\Sigma_\chi) = \emptyset$ for each $W \in \mathscr{W}$; and $\hat{\Sigma}' \cong \hat{\Sigma}_\chi$. \fix{In more detail, let $x \in \hat{\Sigma}_\chi$ with $x \notin \Sigma_\chi \setminus bd(\Sigma_\chi)$. Then $x$ belongs to the closure of some component $r$ of $\hat{\Sigma}_\chi \setminus \Sigma_\chi$. Since $r^*$ is a pole, there exists $W \in \mathscr{W}$ with $r \subseteq W$, by the definition of a cover; and so $x \in W$, and $x \notin bd(W)$ since $bd(W) \subseteq \Sigma_\chi \setminus bd(\Sigma_\chi)$. Hence every point of $\hat{\Sigma}_\chi$ outside $\Sigma_\chi \setminus bd(\Sigma_\chi)$ lies in $W \setminus bd(W)$ for some $W \in \mathscr{W}$. Consequently $\Sigma'$ equals $\hat{\Sigma}_\chi \setminus \bigcup(W \setminus bd(W) : W \in \mathscr{W})$, as a subspace of $\hat{\Sigma}_\chi$. Since the members of $\mathscr{W}$ are mutually disjoint closed discs in the connected surface $\hat{\Sigma}_\chi$, which has empty boundary, it follows that $\Sigma'$ is a connected surface, that its cuffs are the sets $bd(W)$ $(W \in \mathscr{W})$, and that pasting the discs $W \in \mathscr{W}$ onto these cuffs gives $\hat{\Sigma}_\chi$. Thus $\hat{\Sigma}' \cong \hat{\Sigma}_\chi$ and $bd(\Sigma') = \bigcup(bd(W) : W \in \mathscr{W})$.}\fixnote{fix:14.1-surface-added}{The paper gives the reason only for the first of these facts. The argument above is the one used in the formalization.} Let $\Phi'$ be the frame in $\Sigma'$ with
\begin{align*}
U(\Phi') &= bd(\Sigma'), \\
V(\Phi') &= \bigcup(N(W) : W \in \mathscr{W}),
\end{align*}
where the edges of $\Phi'$ are directed arbitrarily, and an edge of $\Phi'$ is designated as short if it includes a member of $S(W)$ for some $W \in \mathscr{W}$, and long otherwise.

Let $(\Gamma, \phi, \lambda, \mathscr{T}) \in \mathscr{D}'$, and let $W \in \mathscr{W}$. For $v \in V(\Gamma) \cap bd(W)$, let the spoke $S_v$ and disc $\Delta_v$ be defined as in 13.5. If $|V(H(W))| > 1$, we define the \textit{hub sector} (of $(\Gamma, \phi, \lambda, \mathscr{T})$ at $W$) to be the rooted hypergraph $(\Gamma \cap ins(H(W)), \pi(W))$. If $|V(H(W))| = 1$ then there is no hub sector. For each $r^* \in (V(\Gamma^*) \setminus V(\Gamma)) \cap bd(W)$, let $r^*$ have neighbours $v_1, v_2$ in $\Gamma^* \cap bd(W)$ and let $A^-$ be the union of $\Gamma \cap \Delta_{r^*}$ and the hypergraph with vertex set $V(S_{v_1} \cap S_{v_2}) \cap V(\Gamma)$ and edge set empty. Let $A$ be the rooted hypergraph $(A^-, \pi)$, where $\pi = \pi(r^*)$ if $r^* \in S(W)$, and otherwise $\pi$ is an arbitrary march with $\bar{\pi} = V(S_{v_1} \cup S_{v_2}) \cap V(\Gamma)$. We call $A$ the $r^*$\textit{-sector} (at $W$). Let $\mathscr{L}(W)$ be the set containing the hub sector at $W$ if there is one and each $r^*$-sector at $W$, for all $r^* \in (V(\Gamma^*) \setminus V(\Gamma)) \cap bd(W)$; and let $\mathscr{L} = \bigcup(\mathscr{L}(W) : W \in \mathscr{W})$.

\hspace*{12pt}(2) \textit{$\mathscr{L}$ $(4n \cdot 5^{2m} + 3)$-isolates $\mathscr{T}$.}

\textbf{Subproof.} We claim first, that each member of $\mathscr{L}$ has order at most $4n \cdot 5^{2m} + 2$. For $|E(H(W))| \le 4\rho(W) + 4$ and so hub sectors have order $\le 2\rho(W) + 2$; for $r^*$ singular, each $r^*$-sector has order
\begin{equation*}
\frac{1}{2}|bd(\Delta_{r^*}) \cap V(\Gamma)| \le \frac{1}{2}(4\rho(W) + 4),
\end{equation*}
by 13.5; and for $r^*$ regular, each $r^*$-sector has order $|V(S_{v_1} \cup S_{v_2}) \cap V(\Gamma)| \le \rho(W) + 1$ where $v_1, v_2$ are the neighbours of $r^*$ in $\Gamma^* \cap bd(W)$. Since each $\rho(W) \le 2n \cdot 5^{2m}$, this proves our claim. Moreover, $\mathscr{L}^- \subseteq \mathscr{T}$ by 13.5. Let $\mathscr{T}'$ be a tangle in $\Gamma$ of order $\ge 4n \cdot 5^{2m} + 3$ such that $(B, A) \in \mathscr{T}'$ for some $(A, B) \in \mathscr{L}^-$ where $(A, B) \in (\mathscr{L}(W))^-$
